% Options for packages loaded elsewhere
\PassOptionsToPackage{unicode}{hyperref}
\PassOptionsToPackage{hyphens}{url}
\documentclass[
]{article}
\usepackage{xcolor}
\usepackage{amsmath,amssymb}
\setcounter{secnumdepth}{-\maxdimen} % remove section numbering
\usepackage{iftex}
\ifPDFTeX
  \usepackage[T1]{fontenc}
  \usepackage[utf8]{inputenc}
  \usepackage{textcomp} % provide euro and other symbols
\else % if luatex or xetex
  \usepackage{unicode-math} % this also loads fontspec
  \defaultfontfeatures{Scale=MatchLowercase}
  \defaultfontfeatures[\rmfamily]{Ligatures=TeX,Scale=1}
\fi
\usepackage{lmodern}
\ifPDFTeX\else
  % xetex/luatex font selection
\fi
% Use upquote if available, for straight quotes in verbatim environments
\IfFileExists{upquote.sty}{\usepackage{upquote}}{}
\IfFileExists{microtype.sty}{% use microtype if available
  \usepackage[]{microtype}
  \UseMicrotypeSet[protrusion]{basicmath} % disable protrusion for tt fonts
}{}
\makeatletter
\@ifundefined{KOMAClassName}{% if non-KOMA class
  \IfFileExists{parskip.sty}{%
    \usepackage{parskip}
  }{% else
    \setlength{\parindent}{0pt}
    \setlength{\parskip}{6pt plus 2pt minus 1pt}}
}{% if KOMA class
  \KOMAoptions{parskip=half}}
\makeatother
\setlength{\emergencystretch}{3em} % prevent overfull lines
\providecommand{\tightlist}{%
  \setlength{\itemsep}{0pt}\setlength{\parskip}{0pt}}
\usepackage{fontspec}
\setmainfont{Courier New}
\setmonofont{Courier New}

\usepackage[letterpaper,left=0.62in,right=0.62in,top=0.52in,bottom=0.55in]{geometry}
\usepackage{xcolor}
\definecolor{resumeaccent}{HTML}{35B4DE}
\definecolor{resumetext}{HTML}{25282A}
\definecolor{resumemuted}{HTML}{6D7377}

\usepackage{microtype}
\usepackage{titlesec}
\usepackage{enumitem}
\usepackage{fancyhdr}
\usepackage{hyperref}

\hypersetup{
  colorlinks=true,
  urlcolor=resumeaccent,
  linkcolor=resumeaccent,
  pdfborder={0 0 0}
}

\color{resumetext}
\setlength{\parindent}{0pt}
\setlength{\parskip}{0.42em}
\linespread{1.03}

\titleformat{\section}
  {\normalfont\fontsize{22}{25}\selectfont\bfseries\color{resumetext}}
  {}{0pt}{}
\titlespacing*{\section}{0pt}{0pt}{0.35em}

\titleformat{\subsection}
  {\normalfont\fontsize{10.5}{12}\selectfont\bfseries\color{resumeaccent}\scshape}
  {}{0pt}{}
  [\vspace{-0.35em}\color{resumetext}\titlerule]
\titlespacing*{\subsection}{0pt}{1.25em}{0.55em}

\titleformat{\subsubsection}
  {\normalfont\fontsize{10.5}{12}\selectfont\bfseries\color{resumetext}}
  {}{0pt}{}
\titlespacing*{\subsubsection}{0pt}{0.8em}{0.15em}

\titleformat{\paragraph}[block]
  {\normalfont\fontsize{9.7}{11}\selectfont\bfseries\color{resumetext}}
  {}{0pt}{}
\titlespacing*{\paragraph}{0pt}{0.55em}{0.05em}

\setlist[itemize]{
  leftmargin=1.25em,
  label=\textbullet,
  itemsep=0.12em,
  topsep=0.18em,
  parsep=0pt,
  partopsep=0pt
}

\fancypagestyle{resume}{
  \fancyhf{}
  \fancyfoot[C]{\color{resumemuted}\fontsize{8}{9}\selectfont\thepage}
  \renewcommand{\headrulewidth}{0pt}
  \renewcommand{\footrulewidth}{0pt}
}
\pagestyle{resume}

\AtBeginDocument{\noindent\color{resumeaccent}\rule{\linewidth}{3pt}\vspace{0.55em}}
\AtBeginDocument{\raggedbottom}

\renewcommand{\tightlist}{%
  \setlength{\itemsep}{0.12em}\setlength{\parskip}{0pt}\setlength{\parsep}{0pt}%
}
\usepackage{bookmark}
\IfFileExists{xurl.sty}{\usepackage{xurl}}{} % add URL line breaks if available
\urlstyle{same}
% fallback for those not using the hyperref driver hyperxmp:
\makeatletter
\@ifundefined{xmpquote}{\newcommand{\xmpquote}[1]{#1}}{}
\makeatother
\hypersetup{
  hidelinks,
  pdfcreator={LaTeX via pandoc}}

\author{}
\date{}

\begin{document}

\section{Anand Sharma}\label{anand-sharma}

\textbf{Architect / Senior Staff Engineer}

Staff Software Engineer, Solutions Architect, DevOps, AI-Pilled

\begin{itemize}
\tightlist
\item
  Austin, TX
\item
  (408) 597-7614
\item
  anand.sharma@gmail.com
\item
  \href{https://www.linkedin.com/in/indrayam/}{LinkedIn}
\end{itemize}

\subsection{Summary}\label{summary}

\href{https://resume.indrayam.com/architect.html}{Software architect}
who still believes the Unix terminal is the most elegant interface ever
invented. I started in graduate school with an SGI workstation running
IRIX; more than two decades later, I bring that same instinct to my work
at Cisco, where I build enterprise DevOps and developer-platform systems
used by 30,000+ engineers. My focus is making complex software delivery
feel simple - from CI/CD, release orchestration, and cloud-native
architecture to AI-assisted tools and multi-agent systems. I care deeply
about the details, but I also understand the importance of zooming out -
distilling the essence of the problem, its context, and its
possibilities so others can see what matters.

I like to build things, evangelize the ideas behind them, and inspire
the people around me to make things that matter.

\subsection{Key Achievements}\label{key-achievements}

\begin{itemize}
\tightlist
\item
  Invited to present at the prestigious Cisco Live 2017 and Cisco Live
  2016 events on advancements in DevOps and Enterprise Class Continuous
  Delivery.
\item
  Led Cisco IT's first Software Developer Experience as a Thought
  Leader, Evangelist, and Solutions Architect.
\item
  Innovated development of Cisco's first centrally managed Continuous
  Delivery Platform.
\item
  Architected the first customer-facing Enterprise App Store for Cisco's
  Cius product.
\item
  Created the first Social Computing Platform for Cisco.
\item
  Oversaw development of the first Cisco API Management Platform,
  featuring a customer-facing API Console.
\end{itemize}

\subsection{Experience}\label{experience}

\subsubsection{Cisco Systems}\label{cisco-systems}

\paragraph{Senior Staff Software
Engineer}\label{senior-staff-software-engineer}

\textbf{Job Title:} Sr.~Staff Software Engineer (Enterprise
DevOps-as-a-Service)\\
\textbf{Dates:} February 2023 - Present\\
\textbf{Location:} Austin, Texas Metropolitan Area - Remote

\begin{itemize}
\tightlist
\item
  Lead efforts to standardize an enterprise-scale PPM tool for
  Operations and Engineering, including conducting a thorough analysis
  of Aha!, Atlassian Jira Align, and ServiceNow.
\item
  Drive widespread adoption of Jira Align across multiple teams
  following a presentation to the COO on key findings.
\item
  Investigate solutions to modernize the stack, reduce technical debt,
  and offer a frictionless experience for the 25K+ Cisco employees
  across Engineering, Customer Support, and Operations, replacing a
  homegrown tool for defect lifecycle management.
\end{itemize}

\paragraph{Staff Software Engineer}\label{staff-software-engineer}

\textbf{Dates:} January 2022 - February 2023\\
\textbf{Location:} Austin, Texas - Remote

\begin{itemize}
\tightlist
\item
  Spearhead the definition, design, documentation, and delivery of
  cutting-edge software delivery (Enterprise DevOps) tools and
  experiences to a team of over 30K software engineers.
\item
  Shape enterprise-wide decision-making by presenting frequent,
  impactful strategic and operational updates related to services and
  service offerings to Senior IT Leadership.
\item
  Provide insightful thought leadership and technical solutions while
  evangelizing solutions to Cisco IT customers.
\item
  Reduced costs by launching the Cisco FCS Software Distribution
  solution using Dev Hub Download to support shipping FCS binaries as
  container images instead of paying Cloud Providers to host FCS
  container images.
\end{itemize}

\paragraph{Senior DevOps Architect}\label{senior-devops-architect}

\textbf{Job Title:} Sr.~Solutions Architect, DevOps\\
\textbf{Dates:} February 2019 - January 2022\\
\textbf{Location:} Raleigh, North Carolina - On-site

\begin{itemize}
\tightlist
\item
  Introduced first-of-its-kind Code Analytics tool to measure software
  delivery performance (quality, security, and release metrics) for over
  approximately 17K unique software, enabling 1K+ engineers to track
  progress on DevSecOps metrics.
\item
  Increased adoption of Cisco IT CI/CD practices for software delivery
  to 90\% across the organization.
\item
  Boosted the percentage of IT applications deployed into production
  with unit tests from 7\% to 66\% and improved unit test coverage for
  applications by 600\%.
\item
  Led event-driven microservices architecture build using a Java
  SpringBoot, Kafka, Redis, and Angular UX stack.
\end{itemize}

\paragraph{DevOps Architect}\label{devops-architect}

\textbf{Job Title:} IT Architect, DevOps\\
\textbf{Dates:} January 2018 - February 2019\\
\textbf{Location:} Raleigh, North Carolina - On-site

\begin{itemize}
\tightlist
\item
  Drove delivery of approximately 50K software promotions to higher
  lifecycle environments by creating Code Release as a single software
  release orchestrator to deliver software using any deployment tool
  anywhere.
\item
  Grew total raw Production Deployments number by 90\% and increased
  median daily Production Deployments by approximately 300\% to
  approximately 100 per day.
\item
  Reduced software promotion times to approximately 10 minutes, which
  improved Deployment Frequency DORA metrics for Cisco IT.
\item
  Led event-driven microservices architecture build using a Java
  SpringBoot, Kafka, Redis, Argo Workflow, and React-based UX stack.
\end{itemize}

\paragraph{DevOps Architect}\label{devops-architect-1}

\textbf{Job Title:} IT Architect, DevOps\\
\textbf{Dates:} March 2016 - January 2018\\
\textbf{Location:} Raleigh, North Carolina - On-site

\begin{itemize}
\tightlist
\item
  Led Cisco IT's first-ever Software Developer Experience Portal as a
  Thought Leader, Evangelist, and Solutions Architect. The Software
  Developer Experience removes friction and contextual work for the
  developer, allowing them to continue to focus on building software.
\item
  Improved the Lead Time for Changes DORA metric by reducing
  self-service requests to approximately 5 minutes, resulting in the
  self-service of approximately 22K CRUD requests across 4K+ unique
  Software Engineers using Code Console in 2018.
\item
  Oversaw event-driven Microservices architecture built using Java
  SpringBoot, MongoDB, Redis, and React-based UX stack.
\end{itemize}

\paragraph{DevOps Architect}\label{devops-architect-2}

\textbf{Job Title:} IT Architect, API Management\\
\textbf{Dates:} February 2014 - March 2016\\
\textbf{Location:} Raleigh, North Carolina - On-site

\begin{itemize}
\tightlist
\item
  Launched Cisco IT's first Continuous Delivery Platform with a mix of
  open-source and commercial tools, resulting in approximately 90\%
  adoption of CI/CD tools and mature software engineering practices to
  deliver software.
\item
  Partnered with Cisco IT Delivery Organizations to successfully deliver
  an enterprise-wide Continuous Delivery solution that scaled across
  approximately 100 IT Services and approximately 2,000 IT Applications.
\item
  Enhanced deployment of IT Apps into production with Unit Tests from
  less than 10\% to 60\%+ by building a homegrown solution to handle
  continuous delivery of Oracle software, addressing a gap in Oracle's
  existing offerings. This solution is currently used by 90\% of Cisco
  teams for 100+ deployments a day.
\item
  Delivered a series of core DevOps tools including GitHub SaaS, GitHub
  Enterprise, CloudBees Jenkins, Artifactory, SonarQube, Spinnaker, and
  CloudBees CDRO.
\end{itemize}

\paragraph{Information Technology System
Architect}\label{information-technology-system-architect}

\textbf{Dates:} July 2010 - February 2014\\
\textbf{Location:} Raleigh, North Carolina - On-site

No detailed description was surfaced in the current LinkedIn profile.

\paragraph{Information Technology Technical
Lead}\label{information-technology-technical-lead}

\textbf{Job Title:} IT Technical Leader, Web 2.0 (Customer Advocacy IT
Team)\\
\textbf{Dates:} September 2005 - July 2010\\
\textbf{Location:} San Jose, California - On-site

No detailed description was surfaced in the current LinkedIn profile.

\paragraph{Information Technology
Engineer}\label{information-technology-engineer}

\textbf{Job Title:} IT Software Engineer (Cisco Advanced Services IT
Team)\\
\textbf{Dates:} August 2001 - September 2005\\
\textbf{Location:} San Jose, California - On-site

\begin{itemize}
\tightlist
\item
  Architected and designed web-based applications for the Advanced
  Services business unit.
\item
  Developed web-based application services on the J2EE platform,
  primarily using open-source MVC frameworks such as Struts and WebWork.
\item
  Continued to act in the capacity of System Administrator of the
  Professional Services Automation application. Member of the core team
  responsible for both the US and Global launch of this application.
\item
  Provided 24x7 systems and application administrative support to all
  AS-IT applications, including ASA (Advanced Services Automation) and
  ASAP (Advanced Services Account Portal).
\end{itemize}

\paragraph{Information Technology
Analyst}\label{information-technology-analyst}

\textbf{Dates:} February 2001 - August 2001

\begin{itemize}
\tightlist
\item
  Responsible for understanding business needs and assisted in authoring
  business requirements documents (BRDs).
\item
  Worked with the Technical Lead in authoring Systems Requirement
  Documents (SRDs).
\item
  Assisted in the formulation of testing and quality assurance scripts
  and plans for applications.
\end{itemize}

\subsubsection{Avaya}\label{avaya}

\paragraph{Web Architect}\label{web-architect}

\textbf{Dates:} March 2000 - February 2001

The New Technology/Architecture group architected the various components
of the e-business infrastructure. It was also responsible for
researching new technologies with the aim of building world-class
e-business services for customers.

\begin{itemize}
\tightlist
\item
  Researched ways to make Avaya's e-business applications WAP-enabled
  and researched new publishing frameworks based on Java and XSLT
  technologies.
\end{itemize}

\subsubsection{Lucent Technologies}\label{lucent-technologies}

\paragraph{IT Software Engineer}\label{it-software-engineer}

\textbf{Dates:} January 1999 - February 2001

\begin{itemize}
\tightlist
\item
  Played a lead role in designing and developing Lucent's first Single
  Sign-on solution for external web applications, leveraging the
  Netegrity SiteMinder solution.
\end{itemize}

\paragraph{Web Technical Leader}\label{web-technical-leader}

\textbf{Dates:} February 1999 - March 2000

Customer Self Support (support.lucent.com) was an extremely busy site.
As the lead developer of the first SSO pilot, took on the added
responsibility of leading the effort to migrate the entire application
into an NES presentation layer, NAS (Netscape Application Server)
business-logic layer, and Oracle 7.3.x data layer framework.

\begin{itemize}
\tightlist
\item
  Led a team of three developers and rewrote all server-side programs
  into Java applogics.
\item
  Administered the development NES/NAS/Oracle environment for the CSS
  project.
\end{itemize}

\subsubsection{Arizona State University}\label{arizona-state-university}

\paragraph{IS Lab Lead}\label{is-lab-lead}

\textbf{Dates:} May 1996 - December 1998

The IS Lab catered to the technology needs of Arizona State instructors
and students. It focused on helping the community embrace new and
upcoming technologies in education.

\begin{itemize}
\tightlist
\item
  Worked on projects involving web applications, proofs of concept for
  new technologies and products, and helping IS Lab customers get
  introduced to them.
\end{itemize}

\subsection{Education}\label{education}

\subsubsection{Master of Computer Science, Software
Engineering}\label{master-of-computer-science-software-engineering}

Arizona State University - Tempe, AZ

\subsubsection{Master of Science, Industrial
Engineering}\label{master-of-science-industrial-engineering}

Arizona State University - Tempe, AZ

\subsection{Awards and Presentations}\label{awards-and-presentations}

\subsubsection{2018 Gold Stevie Winner}\label{gold-stevie-winner}

Enterprise Software Release Digitization, Secure DevOps

\subsubsection{Cisco Live 2016}\label{cisco-live-2016}

Panel discussion: ``Addressing the What's, Why's and How's in the New
World of DevOps''

\subsubsection{Cisco Live 2017}\label{cisco-live-2017}

Presentation: ``Enterprise Class Continuous Delivery''

\subsubsection{Cloud Identity Summit 2012 - Vail,
CO}\label{cloud-identity-summit-2012---vail-co}

Presentation:
\href{https://www.slideshare.net/slideshow/who-says-elephant-cant-dance/13698356?utm_source=clipboard_share_button&utm_campaign=slideshare_make_sharing_viral_v2&utm_variation=variant&utm_medium=share}{``Who
Says Elephant Can't Dance? Securely Externalizing APIs @ Cisco''} hosted
by Ping Identity

\end{document}
